% Programming with InitechBase, Chapter 3: Memory Variables (AN-16; annex
% Sec 8.2 P3).  Panels replayed (examples/p3-*.dot); sources cmd/flow.c
% (STORE, =, the memvar table, 256 variables, names of 10), cmd/proc.c
% (PUBLIC, PRIVATE, scope), core/fn_builtins.c TYPE.
\chapter{Memory Variables}
\Gates{test-pubs-examples, test-interp-flow, test-interp-proc}
A program needs places to keep values that are not fields of a database: a
total, a count, an answer, a date.  These are memory variables.  This chapter
describes how they are created, what types they have, and how long they live.

\begin{ChapterCovers}
\item Creating memory variables
\item Types
\item Totals and counts
\item Public and private memory variables
\item Releasing memory variables
\end{ChapterCovers}

\ATEHead{Preparing for This Chapter}
Start InitechBase with a fresh copy of the sample diskette.

\ATEHead{Creating Memory Variables}
A memory variable is created when a value is first stored in it, by
\cmd{STORE} or by the equals sign:

\SyntaxLine{\kw{STORE} \xp{exp} \kw{TO} \xp{name list}}
\SyntaxLine{\xp{name} \kw{=} \xp{exp}}

\cmd{STORE} can give several memory variables the same value at once; the
programs of the system use it to set their totals to zero.  A name has up to 10
letters, digits and underscores and begins with a letter.  The system begins
every name with M: \texttt{MTOTAL}, \texttt{MCOUNT}, \texttt{MINVNO}.  A name that
is also the name of a field would always stand for the field.

\ATEHead{Types}
A memory variable takes the type of the value stored in it: character,
numeric, date or logical.  The function \fn{TYPE()} tells you the type of a
value:
\DotExample{examples/p3-types}
Storing a value of another type in the memory variable changes its type.

\ATEHead{Totals and Counts}
A total is a memory variable that starts at zero and grows by an amount for
each record; a count grows by one.  \file{INVTOTAL.PRG}, listed in Chapter~2,
keeps two totals:
\DotExample{examples/p3-invtotal}

\ATEHead{Public and Private Memory Variables}
The memory variables created by the main program of a file remain after it
ends, as if you had created them at the dot prompt.  Those created inside a
procedure (Chapter~4) are private to it: they are released when the procedure
returns.

\cmd{PUBLIC} declares memory variables that every part of every program can use
and that remain until they are released.  \file{INVTOTAL.PRG} declares its two
totals \cmd{PUBLIC}, so that the month-end journal entry can be made from them
after the program has run.  Declare a memory variable \cmd{PUBLIC} before you
store a value in it.

\cmd{PRIVATE} declares memory variables that belong to the procedure that
declares them.  A memory variable of the same name created by the caller is
hidden while the procedure runs and is unchanged afterwards.  A procedure that
uses a name its caller may also use, such as \texttt{MCOUNT}, declares it
\cmd{PRIVATE}.

\ATEHead{Releasing Memory Variables}
\cmd{RELEASE} removes memory variables, and \cmd{RELEASE ALL} removes all of them.
InitechBase holds up to 256 memory variables; a program that creates many in a
loop releases those it no longer needs.

\ATEHead{Where Do You Go From Here?}
Chapter~4 divides a program into procedures and passes values to them.

\ATEHead{Summary}
\begin{enumerate}
\item \cmd{STORE} and \texttt{=} create memory variables and change their values.
\item A memory variable takes the type of its value; \fn{TYPE()} tells it.
\item A total starts at zero and grows for each record.
\item \cmd{PUBLIC} memory variables remain; \cmd{PRIVATE} ones belong to their
  procedure.
\item \cmd{RELEASE} removes memory variables.
\end{enumerate}
